\chapter{Situations and Applications}\label{ch:situations}

A \emph{situation} is a requirement (an interface, a behavior, a budget) and a \emph{solution} is a molecule that meets it. The decision matrices of the software standard are exactly the formal optimization over the set of solutions; the two theorems of this chapter make that set finite and decidable (\autoref{thm:51}) and make zero-copy elimination sound (\autoref{thm:52}).

\paragraph{Situations.}
Fix the finite type universe of the dataplane and the cost model of \autoref{ch:enrich}, which assigns every molecule a cost vector $\mathrm{cost} (M) = (t (M), m (M), c (M))$(time, memory, cache misses) additive under composition (\autoref{thm:25}) and monotone under refinement (\autoref{thm:28}).

A \emph{situation} is a triple $\sigma = (A \to B,\, \Phi,\, \bar c)$:
\begin{itemize}
\item $A \to B$: the interface, with finite input type $A$ and finite output type $B$ (\emph{bounded types});
\item $\Phi$: a behavioral specification, a decidable predicate on molecules $A \to B$ invariant under behavioral equivalence: e.g.\ ``$M$ realizes the regular language $L$'' (decidable by minimization, \autoref{ch:coalgebra});
\item $\bar c = (c_T, c_M, c_C)$: a cost budget vector, where the memory component counts the states of the canonical presentation, $m (M) = |S_M|$.
\end{itemize}

A molecule $M : A \to B$ is a \emph{solution} of $\sigma$, written $M \models \sigma$, when $M \vdash \Phi$ and $\mathrm{cost} (M) \leq \bar c$ (componentwise); the feasible set is $F (\sigma) = \{ M : A \to B \mid M \models \sigma \}$. A \emph{decision matrix} for $\sigma$ is the computation
\[
M^{*} = \operatorname*{arg\,min}_{M \in F (\sigma)} \mathrm{cost} (M),
\]
returning \emph{infeasible} when $F (\sigma) = \emptyset$. Under the refinement preorder of \autoref{ch:enrich} the feasible set is a down-set: refinement preserves behavior and is cost-monotone (\autoref{thm:28}), so $M' \sqsubseteq M \models \sigma$ implies $M' \models \sigma$; the decision matrix therefore optimizes downward through the refinement poset, and by maximal refinement (\autoref{thm:29}) its output is the most-refined feasible molecule.

\setcatalog{51}
\begin{theorem}[Situation solvability (\autoref{thm:51})]\label{thm:51}
Let $\sigma = (A \to B, \Phi, \bar c)$ be a situation with finite types $A, B$, a decidable specification $\Phi$, and a memory budget $c_M$ that bounds state size via $m (M) = |S_M|$. Then:
\begin{enumerate}
\item the feasible set $F (\sigma)$ is finite;
\item membership in $F (\sigma)$ is decidable: for every molecule $M : A \to B$, both $M \vdash \Phi$ and $\mathrm{cost} (M) \leq \bar c$ are decidable predicates;
\item the decision-matrix computation, the minimization over $F (\sigma)$, terminates.
\end{enumerate}
\end{theorem}

\begin{proof}
(i) A solution satisfies $m (M) \leq c_M$, hence $|S_M| \leq K$ with $K = \lfloor c_M \rfloor$. Fix a state set $S = \{1, \ldots, K\}$: a Mealy machine $A \to B$ with state space $S$ is determined by its initial state $s_0 \in S$ and its transition $\delta : S \times A \to B \times S$, so there are at most $K \cdot (|B| \cdot K)^{K \cdot |A|}$ such machines, a finite number. The molecules of Mol are the state-isomorphism classes of such machines, and a quotient of a finite set is finite; since every solution has at most $K$ states, $F (\sigma)$ is finite.

(ii) For a fixed molecule $M$, presented finitely as a Mealy machine, the predicate $M \vdash \Phi$ is decidable by hypothesis on $\Phi$. The cost $\mathrm{cost} (M)$ is computed from the declared constant cost vectors of the atoms of $M$ by the additive rules (\autoref{thm:25}): composition adds componentwise and the tensor combines components; the result is a vector of rationals, and comparing it with $\bar c$ is decidable. Hence membership in $F (\sigma)$ is decidable.

(iii) By (i) the candidate set is finite and by (ii) membership is decidable, so enumerating $F (\sigma)$ and taking the minimum (or reporting infeasibility when $F (\sigma)$ is empty) terminates.
\end{proof}

\paragraph{Well-formed domains and zero copy.}
Let $B$ be the byte-stream type and $T$ a structured type. A parser is a molecule $\mathrm{parse} : B \to T$ and a serializer is a molecule $\mathrm{serialize} : T \to B$. A subset $D \subseteq T$ is \emph{well-formed} when the roundtrip property holds on it: $\mathrm{parse} \circ \mathrm{serialize} = \mathrm{id}$ on $D$, that is, $\mathrm{parse} (\mathrm{serialize} (t)) = t$ for every $t \in D$. The \emph{well-formed byte domain} is $W = \mathrm{serialize} (D) \subseteq B$.

\setcatalog{52}
\begin{theorem}[Zero-copy roundtrip (\autoref{thm:52})]\label{thm:52}
Let $\mathrm{parse} : B \to T$ and $\mathrm{serialize} : T \to B$ be molecules with $\mathrm{parse} \circ \mathrm{serialize} = \mathrm{id}$ on a well-formed domain $D \subseteq T$, and let $W = \mathrm{serialize} (D)$. Then:
\begin{enumerate}
\item $\mathrm{serialize} \circ \mathrm{parse} = \mathrm{id}$ on $W$;
\item for every molecule $X : T \to T$ with $X (D) \subseteq D$ and $X = \mathrm{id}$ on $D$, the roundtrip segment $\mathrm{serialize} \circ X \circ \mathrm{parse} : B \to B$ is the identity on $W$;
\item the elimination is sound: if $X \neq \mathrm{id}$ on $D$, then $\mathrm{serialize} \circ X \circ \mathrm{parse} \neq \mathrm{id}$ on $W$, so replacing the segment by a pass-through would change behavior on well-formed traffic.
\end{enumerate}
\end{theorem}

\begin{proof}
(i) For $b \in W$, write $b = \mathrm{serialize} (t)$ with $t \in D$. The roundtrip property gives $\mathrm{parse} (b) = t$, hence
\[
(\mathrm{serialize} \circ \mathrm{parse})(b) = \mathrm{serialize} (\mathrm{parse} (b)) = \mathrm{serialize} (t) = b .
\]

(ii) For $b \in W$, write $b = \mathrm{serialize} (t)$ with $t \in D$. Since $\mathrm{parse} (b) = t$ and $X = \mathrm{id}$ on $D$,
\[
(\mathrm{serialize} \circ X \circ \mathrm{parse})(b) = \mathrm{serialize}\bigl (X (\mathrm{parse} (b))\bigr) = \mathrm{serialize} (X (t)) = \mathrm{serialize} (t) = b .
\]

(iii) The roundtrip property makes $\mathrm{serialize}$ injective on $D$: if $t, t' \in D$ and $\mathrm{serialize} (t) = \mathrm{serialize} (t')$, then $t = \mathrm{parse} (\mathrm{serialize} (t)) = \mathrm{parse} (\mathrm{serialize} (t')) = t'$. If $X \neq \mathrm{id}$ on $D$, take $t \in D$ with $X (t) \neq t$; since $X (D) \subseteq D$, both $t$ and $X (t)$ lie in $D$, and injectivity gives $\mathrm{serialize} (X (t)) \neq \mathrm{serialize} (t)$. For $b = \mathrm{serialize} (t) \in W$: $(\mathrm{serialize} \circ X \circ \mathrm{parse})(b) = \mathrm{serialize} (X (t)) \neq \mathrm{serialize} (t) = b$, so the segment differs from the identity on well-formed input and eliminating it would change behavior.
\end{proof}

\paragraph{Remark.}
This is the license for decision matrix D-6 (zero-copy vs.\ copy): when $X$ is a no-op on the well-formed domain, \autoref{thm:52} guarantees that a proxy's parse--process--serialize segment may be replaced by a pass-through that forwards the buffer (no parse, no serialize, no allocation) while behavior on well-formed traffic is preserved; when $X$ moves within $D$, part (iii) shows the copy is semantically necessary.

\paragraph{Precomputed dispatch.}
By \autoref{thm:51} the decision matrix is a finite computation at design time, and its output is a finite table of (discriminant, molecule) pairs. At runtime, dispatch is a lookup keyed by a dense integer discriminant: the precomputed lookup table of \autoref{ch:free} (\autoref{thm:24}): no dynamic dispatch, no search, a single indexed jump. The decision-matrix output is thus exactly the runtime dispatch table.

\paragraph{Canonical situations.}
The software standard treats four situations as worked recipes; each is a special case of the formal setup above.
\begin{itemize}
\item \textbf{Web server (HTTP over TCP).} $B = $ TCP byte stream, $T = $ HTTP request; $\Phi$: RFC-conformant request framing; $\bar c$: bounded per-request latency and memory. The solution is a pipeline of linear atoms (\autoref{thm:50}): a ring-fed parser (\autoref{thm:48}), a framing atom, a dispatch atom (\autoref{thm:51} lookup); the batch size is set by the amortization curve of \autoref{thm:47} (matrix D-4), and for proxies \autoref{thm:52} makes a no-op relay a pass-through.
\item \textbf{DNS (UDP and TCP).} Two transport modes, selected by matrix D-3; DNS messages are small, so batches are bounded by the ring invariant and the amortization law (\autoref{thm:47}, \autoref{thm:48}), and the response molecule is precomputed by \autoref{thm:51}.
\item \textbf{FTP (control and data).} Two coupled pipelines: a session state machine over the control channel, a bulk transfer path over the data channel, joined by a shared session-state atom; the memory budget is split between the two rings (\autoref{thm:48}).
\item \textbf{Custom protocol.} A fixed binary protocol gives a finite interface: \autoref{thm:51} enumerates the feasible set, \autoref{thm:24} precomputes the dispatch table, and the parser/serializer pair is inverted on its well-formed domain (\autoref{thm:52}, zero copy) or kept with copy.
\end{itemize}

The decision matrices of the standard and the results that govern them are collected in \autoref{tab:decision-matrices}.

\begin{table}[h]
\centering
\small
\begin{tabular}{@{}lll@{}}
\toprule
Matrix & Decision & Governing result \\
\midrule
D-1  & Ring/buffer sizing (L3-aware)        & \autoref{thm:48} \\
D-2  & SoA vs.\ AoS layout                  & cost model, \autoref{ch:enrich} \\
D-3  & Protocol selection (UDP/TCP/SCTP)    & \autoref{thm:51} (feasibility) \\
D-4  & Batch size                           & \autoref{thm:47} \\
D-5  & Polling: busy-poll vs.\ SQPOLL vs.\ AF\_XDP & \autoref{thm:47} \\
D-6  & Zero-copy vs.\ copy                  & \autoref{thm:52} \\
D-7  & Dispatch: monomorphization vs.\ indirect & \autoref{thm:24} \\
D-8  & SIMD width                           & cost model, \autoref{ch:enrich} \\
D-9  & Lock-free vs.\ locking               & cost model, \autoref{ch:enrich} \\
D-10 & Hot/cold field split                 & cost model, \autoref{ch:enrich} \\
D-11 & Allocation policy                    & \autoref{thm:50} \\
D-12 & Plugin placement                     & \autoref{thm:51} (feasibility) \\
\bottomrule
\end{tabular}
\caption{The decision matrices of the standard and the results that govern them.}
\label{tab:decision-matrices}
\end{table}

In every case a solution is a molecule $M$ with $M \vdash \Phi$ and $\mathrm{cost} (M) \leq \bar c$, found by the terminating decision-matrix computation of \autoref{thm:51} and refined to zero copy by \autoref{thm:52}. The situations of the standard mirror this chapter by construction, so each policy cites one of the results \autoref{thm:46}--\autoref{thm:52} instead of restating a proof.
